\section{Reproducibility and audit details}

The supplementary release records the encoded-byte hashes used for duplicate
audits, every excluded item, guarded-block coordinates, all retained crop
rectangles, and the exhaustive cross-split intersection result. Seed-level
checkpoints and raw score arrays are retained so that AUROC, AUPRC, FPR95,
threshold-transfer, and bootstrap calculations can be independently
recomputed.

Compressed per-example arrays contain labels, predictions and all synthetic
audit scores for every corpus/seed. The paired measured-GSD arrays additionally
contain UAS and satellite predictions, confidence, received consistency, site,
event, UAS block, satellite block and joint overlap-cluster identifiers.
Portable indexes omit local filesystem paths. The reproduction command verifies
their hashes, recomputes headline metrics and the preregistered W13
adjudication, regenerates publication artefacts and clean-builds the paper.

Failed and superseded branches remain part of the release. In particular, the
initial unpaired measured-GSD criterion, the overlapping-crop analysis, the
reference-dependent operational interpretation, and the failed
received-image-only W8/W9 experiment are not deleted when corrected evidence is
generated.

The first implementation applied an all-seed W13 rule although the frozen
criterion was mean-based. The corrected adjudication records W13 as passed and
the all-seed requirement as a failed post-hoc robustness sensitivity. Other
independent-review sensitivities, including anchor-matched finer-view scoring
and joint UAS/satellite overlap clustering, are explicitly labelled post hoc.

The structured information-set scoping audit is released as
\texttt{information\_set\_literature\_audit.json}. It records six search
queries, 12 included primary works, task/protocol classifications, and the
test-time information available to each method. The audit found no exact
published unavailable-fine-reference deployment comparison in this bounded
set; it is not presented as an exhaustive prevalence estimate.
